% GENERATED FILE. Edit canonical lessons, then run npm run generate:books.
% canonical-source-hash: fnv1a64:c5ab6d3dd7128dc9
% canonical-lessons: PT-C321-carnaval, PT-C321-desfile, PT-C321-mascara, PT-C321-disfarce, PT-C321-disfarcar

\chapter{No Carnaval}
\label{ch:pt-no-carnaval}
\hlchaptermodality{321}

\begin{chapteropening}
\textbf{By the end of this chapter:} \emph{I can describe carnival: the parade, masks and disguises.}

Complete the last lesson of chapter 321: I can describe carnival: the parade, masks and disguises.
\end{chapteropening}

\section[o carnaval]{o carnaval — the carnival}

\label{lesson:PT-C321-carnaval}



\begin{quote}
Before the new one: say the Portuguese for the pilgrim, then the Portuguese for the pilgrimage.
\end{quote}

\subsection*{You'll want to know: o carnaval}
\textbf{o carnaval} — \textquotedblleft{}the carnival\textquotedblright{}.

\begin{quote}
No carnaval, as crianças vão mascaradas para a escola.
\end{quote}

\begin{quote}
\emph{At carnival the children go to school in fancy dress.}
\end{quote}

The days before Lent; often written with a capital, \emph{o Carnaval}.

\begin{cousinweb}
\emph{carnaval} came from Italian \emph{carnevale} (EtymDB), from Medieval Latin \emph{carnelevale} \textquotedblleft{}the putting away of meat\textquotedblright{}, built on \emph{carō} \textquotedblleft{}flesh\textquotedblright{} and \emph{levāre} \textquotedblleft{}to lift\textquotedblright{}.
\end{cousinweb}

\begin{grammarlens}[title={What you've built}]
One more word for telling, explaining and describing at B1.
\end{grammarlens}

\subsection*{Your turn}
\begin{itemize}
\raggedright
  \item \emph{Say it:} \emph{o carnaval}
  \item \emph{Say it:} \emph{o carnaval}, once more
  \item \emph{Say it:} a sentence of your own with \emph{o carnaval}
  \item \emph{Recall:} say the Portuguese for the pilgrim, then the Portuguese for the pilgrimage, then say \emph{o carnaval} again
\end{itemize}

\begin{tcolorbox}[breakable,colback=teal!4,colframe=teal!35!black,title={Before you move on}]
What does o carnaval mean? (\textquotedblleft{}The carnival\textquotedblright{}.) Say it once more.
\end{tcolorbox}
\section[o desfile]{o desfile — the parade}

\label{lesson:PT-C321-desfile}



\begin{quote}
Before the new one: say the Portuguese for the pilgrimage, then the Portuguese for the carnival.
\end{quote}

\subsection*{You'll want to know: o desfile}
\textbf{o desfile} — \textquotedblleft{}the parade\textquotedblright{}.

\begin{quote}
O desfile passou pela rua principal.
\end{quote}

\begin{quote}
\emph{The parade went down the main street.}
\end{quote}

\emph{Um desfile de moda} — a fashion show.

\begin{cousinweb}
\emph{desfile} is built on a verb that came from French \emph{défiler} \textquotedblleft{}to march in file\textquotedblright{}, from \emph{fil} \textquotedblleft{}a thread\textquotedblright{}, which goes back to Latin \emph{fīlum} \textquotedblleft{}a thread\textquotedblright{} (Lewis \& Short).
\end{cousinweb}

\begin{grammarlens}[title={What you've built}]
One more word for telling, explaining and describing at B1.
\end{grammarlens}

\subsection*{Your turn}
\begin{itemize}
\raggedright
  \item \emph{Say it:} \emph{o desfile}
  \item \emph{Say it:} \emph{o desfile}, once more
  \item \emph{Say it:} a sentence of your own with \emph{o desfile}
  \item \emph{Recall:} say the Portuguese for the pilgrimage, then the Portuguese for the carnival, then say \emph{o desfile} again
\end{itemize}

\begin{tcolorbox}[breakable,colback=teal!4,colframe=teal!35!black,title={Before you move on}]
What does o desfile mean? (\textquotedblleft{}The parade\textquotedblright{}.) Say it once more.
\end{tcolorbox}
\section[a máscara]{a máscara — the mask}

\label{lesson:PT-C321-mascara}



\begin{quote}
Before the new one: say the Portuguese for the carnival, then the Portuguese for the parade.
\end{quote}

\subsection*{You'll want to know: a máscara}
\textbf{a máscara} — \textquotedblleft{}the mask\textquotedblright{}.

\begin{quote}
A criança usava uma máscara de leão.
\end{quote}

\begin{quote}
\emph{The child wore a lion mask.}
\end{quote}

Also a face mask for health: \emph{usar máscara}.

\begin{cousinweb}
\emph{máscara} came from Italian \emph{maschera} \textquotedblleft{}a mask, a disguise\textquotedblright{} (EtymDB), from Medieval Latin \emph{masca}.
\end{cousinweb}

\begin{grammarlens}[title={What you've built}]
One more word for telling, explaining and describing at B1.
\end{grammarlens}

\subsection*{Your turn}
\begin{itemize}
\raggedright
  \item \emph{Say it:} \emph{a máscara}
  \item \emph{Say it:} \emph{a máscara}, once more
  \item \emph{Say it:} a sentence of your own with \emph{a máscara}
  \item \emph{Recall:} say the Portuguese for the carnival, then the Portuguese for the parade, then say \emph{a máscara} again
\end{itemize}

\begin{tcolorbox}[breakable,colback=teal!4,colframe=teal!35!black,title={Before you move on}]
What does a máscara mean? (\textquotedblleft{}The mask\textquotedblright{}.) Say it once more.
\end{tcolorbox}
\section[o disfarce]{o disfarce — the disguise; a fancy-dress outfit}

\label{lesson:PT-C321-disfarce}



\begin{quote}
Before the new one: say the Portuguese for the parade, then the Portuguese for the mask.
\end{quote}

\subsection*{You'll want to know: o disfarce}
\textbf{o disfarce} — \textquotedblleft{}the disguise; a fancy-dress outfit\textquotedblright{}.

\begin{quote}
Fiz o meu disfarce com papel e cola.
\end{quote}

\begin{quote}
\emph{I made my fancy-dress outfit out of paper and glue.}
\end{quote}

For carnival, people also say \emph{a fantasia} in Brazil.

\begin{cousinweb}
\emph{disfarce} is built on \emph{disfarçar} \textquotedblleft{}to disguise\textquotedblright{}; the origin of the verb is not settled, and EtymDB has no entry.
\end{cousinweb}

\begin{grammarlens}[title={What you've built}]
One more word for telling, explaining and describing at B1.
\end{grammarlens}

\subsection*{Your turn}
\begin{itemize}
\raggedright
  \item \emph{Say it:} \emph{o disfarce}
  \item \emph{Say it:} \emph{o disfarce}, once more
  \item \emph{Say it:} a sentence of your own with \emph{o disfarce}
  \item \emph{Recall:} say the Portuguese for the parade, then the Portuguese for the mask, then say \emph{o disfarce} again
\end{itemize}

\begin{tcolorbox}[breakable,colback=teal!4,colframe=teal!35!black,title={Before you move on}]
What does o disfarce mean? (\textquotedblleft{}The disguise; a fancy-dress outfit\textquotedblright{}.) Say it once more.
\end{tcolorbox}
\section[disfarçar]{disfarçar — to disguise; to hide (a feeling)}

\label{lesson:PT-C321-disfarcar}



\begin{quote}
Before the new one: say the Portuguese for the mask, then the Portuguese for the disguise.
\end{quote}

\subsection*{You'll want to know: disfarçar}
\textbf{disfarçar} — \textquotedblleft{}to disguise; to hide (a feeling)\textquotedblright{}.

\begin{quote}
Ele disfarçou o medo com um sorriso.
\end{quote}

\begin{quote}
\emph{He hid his fear behind a smile.}
\end{quote}

\emph{Disfarçar-se de} — to dress up as: \emph{disfarçou-se de pirata}.

\begin{cousinweb}
\emph{disfarçar} is usually linked to \emph{des-} and an old word for a face or a farce, but the origin is not settled; EtymDB has no entry.
\end{cousinweb}

\begin{grammarlens}[title={What you've built}]
One more word for telling, explaining and describing at B1.
\end{grammarlens}

\subsection*{Your turn}
\begin{itemize}
\raggedright
  \item \emph{Say it:} \emph{disfarçar}
  \item \emph{Say it:} \emph{disfarçar}, once more
  \item \emph{Say it:} a sentence of your own with \emph{disfarçar}
  \item \emph{Recall:} say the Portuguese for the mask, then the Portuguese for the disguise, then say \emph{disfarçar} again
\end{itemize}

\begin{tcolorbox}[breakable,colback=teal!4,colframe=teal!35!black,title={Before you move on}]
What does disfarçar mean? (\textquotedblleft{}To disguise; to hide (a feeling)\textquotedblright{}.) Say it once more.
\end{tcolorbox}
